\begin{frame}{The BTW sandpile}
    \textbf{Definition:}

    \begin{itemize}
        \item A $d$-dimensional hypercubic grid $R\subset\bZ^d$ of linear size $N$.
        \item A function $z(r):\; \bZ^d \to \bZ$. (The local \enquote{slope} or \enquote{energy}.)
        \item A threshold $z_c$, s.t. for $z(r)>z_c$ the system is unstable.
        \item A relaxation function $z_\tau(r)\mapsto z_{\tau+1}(r)$.
    \end{itemize}
\end{frame}


\begin{frame}{The relaxation function}
    For all $r\in R$ with $z_\tau(r) > z_c$, i.e. at every unstable lattice site:
    \begin{align*}
        z_\tau(r) &\to z_\tau(r) - 2d, \\
        z_\tau(r\pm e_i) &\to z_\tau(r\pm e_i) + 1 \qq{for} i=1,\ldots,d.
    \end{align*}
    Imagine sandpile \enquote{toppling} onto its neighbors.
\end{frame}


\begin{frame}{An avalanche}
    \centering
    %\animategraphics[loop,width=\textwidth]{1}{../../figures/heatmaps/relaxation_example/frame-}{0}{7}
    % TODO: Insert animation of sandpile that starts of in a very unstable state and topples toward stable state (without further perturbation)
\end{frame}


\begin{frame}{Boundary conditions}

    \textbf{Open:}
    \begin{align*}
        z_\tau(r) &\to z_\tau(r) - 2d + (\text{number of }i\text{ with }r_i=N), \\
        z_\tau(r+e_i) &\to z_\tau(r+e_i) + 1 \qq{if} r_i\neq N, \\
        z_\tau(r-e_i) &\to z_\tau(r-e_i) + 1 \qq{for} i=1,\ldots,d, \\
        \text{and} z_\tau(r) &= 0 \qq{if} r_i=1 \text{ for any i}.
    \end{align*} \pause

    \textbf{Closed:}
    \begin{equation*}
        z_\tau(r) = 0 \qq{if} r_i=1\text{ or } r_i=N \text{ for any i}.
    \end{equation*}

\end{frame}


\begin{frame}{Driving the system}
    \textbf{Conservative perturbation:} Keeps the total energy constant.
    \begin{align*}
        z(r_0) &\to z(r_0) + d, \\
        z(r_0-e_i) &\to z(r_0-e_i) - 1 \qq{for} i=1,\ldots,d
    \end{align*}
    Local rearrangement of the sandpile. \pause

    \textbf{Nonconservative perturbation:} Locally increases the energy.
    \begin{equation*}
        z(r_0) \to z(r_0) + 1
    \end{equation*}
    Randomly drop a grain of sand.
\end{frame}


\begin{frame}{Driving the system}

    % TODO: Timelapse animation of empty system being driven until some avalanches happen.
    
    \centering
    %\animategraphics[width=\textwidth]{15}{../../figures/heatmaps/burn_in/frame-}{0}{499}
\end{frame}




\begin{frame}{Reaching criticality}

    \begin{itemize}
        \item Perturbation: Introduces energy to the system
        \item Relaxation: Dissipates energy from the system at the boundaries
        \item In the stationary state: \begin{equation*}
            \expval{z} \to \overline{\expval{z}} = \lim_{T\to\infty}\frac{1}{T}\int_{0}^{T}\expval{z}(\tau)\dd{\tau}.
        \end{equation*}
    \end{itemize}

    In the critical state, the total (spatial average) energy fluctuates around a temporal average energy.

\end{frame}


\begin{frame}{The critical state}

    % TODO: Animation of a system already in stationary state, exhibiting avalanches of various sizes
    \centering
    %\animategraphics[width=\textwidth]{15}{../../figures/heatmaps/stationary/frame-}{0}{499}
\end{frame}


\begin{frame}{What do we measure?}

    \textbf{Avalanche variables:}
    \begin{itemize}
        \item Avalanche lifetime $t$.
        \item Total size $s$. \\ (Cumulative sum of all toppling events.)
        \item Spatial linear size $l$. \\ (Maximum linear extend in any dimension.)
    \end{itemize}

\end{frame}


\begin{frame}{What do we expect?}

    \textbf{Power laws:}
    \begin{align*}
        P(T=t) &\propto t^{1-\alpha}, \\
        P(S=s) &\propto s^{1-\tau}, \\
        P(L=l) &\propto l^{1-\lambda}.
    \end{align*}

\end{frame}
